\section*{第 \ref{ch:b1:stereochemistry-in-depth} 章——立体化学深入}

\begin{solution}{exo:b1:stereochemistry-in-depth:1}
\omterm{def:b1:stereochemistry-in-depth:isomers}{构造异构体}；\omterm{def:b1:stereochemistry-in-depth:enantiomers}{对映异构体}；\omterm{def:b1:stereochemistry-in-depth:enantiomers}{非对映异构体}（即不互为
镜像的\omterm{def:b1:stereochemistry-in-depth:isomers}{立体异构体}）；同一化合物的两种\omterm{def:b1:stereochemistry-in-depth:isomers}{构象}；\omterm{def:b1:stereochemistry-in-depth:enantiomers}{非对映异构体}
（\cip{R,S} 为\omterm{def:b1:stereochemistry-in-depth:enantiomers}{内消旋化合物}）。
\end{solution}

\begin{solution}{exo:b1:stereochemistry-in-depth:2}
\ce{-OH} $>$ \ce{-NH2} $>$ \ce{-CH3} $>$ \ce{-H} (O, N, C, H).
\ce{-COOH} (O,O,O) $>$ \ce{-CHO} (O,O,H) $>$ \ce{-CH2OH} (O,H,H) $>$
\ce{-CH(CH3)2} (C,C,H). \ce{-Br} $>$ \ce{-Cl} $>$ \ce{-CCl3} (Cl,Cl,Cl)
$>$ \ce{-CH2Cl}（Cl,H,H）：直接相连的原子先决定次序，之后才看其邻接原子。
\end{solution}

\begin{solution}{exo:b1:stereochemistry-in-depth:3}
2,3-二氯丁烷：两个等价的\omterm{def:b1:stereochemistry-in-depth:stereocentre}{手性}碳原子，\cip{R,R}、\cip{S,S}和
\omterm{def:b1:stereochemistry-in-depth:enantiomers}{内消旋化合物} \cip{R,S}：共三种。2-溴-3-氯丁烷：两个不同的\omterm{def:b1:stereochemistry-in-depth:stereocentre}{手性}碳原子，
四种\omterm{def:b1:stereochemistry-in-depth:isomers}{立体异构体}，无\omterm{def:b1:stereochemistry-in-depth:enantiomers}{内消旋化合物}。戊烷-2,4-二醇：共三种，\cip{2R,4S}
为\omterm{def:b1:stereochemistry-in-depth:enantiomers}{内消旋化合物}。
\end{solution}

\begin{solution}{exo:b1:stereochemistry-in-depth:4}
在 C5 上：异丙烯基（C,C,C，双键计两次）$>$ C6 $>$ C4
$>$ H。C6 与 C4 都带有（C,H,H）；再往下一步，C6 通向
羰基碳（O,O,C），C4 通向 C3（C,C,H），故 C6 优先。将异丙烯基
朝前、H 朝后，次序异丙烯基 $\to$ C6 $\to$ C4
按图示呈顺时针，\cip{R}，即留兰香味的香芹酮；虚线者为
\cip{S}。
\end{solution}

\begin{solution}{exo:b1:stereochemistry-in-depth:5}
$\theta = 0$：甲基处于\omterm{def:b1:stereochemistry-in-depth:conformations}{重叠式构象}，\qty{16.5}{kJ/mol}；$60^\circ$：\omterm{def:b1:stereochemistry-in-depth:conformations}{邻位交叉构象}，约
\qty{2.9}{kJ/mol}（最小值 2.8，位于 $62^\circ$ 处）；$120^\circ$：甲基与 H 重叠，
\qty{15.2}{kJ/mol}；$180^\circ$：\omterm{def:b1:stereochemistry-in-depth:conformations}{对位交叉构象}，能量为 0。等能量时，两个\omterm{def:b1:stereochemistry-in-depth:conformations}{邻位交叉构象}
对应一个\omterm{def:b1:stereochemistry-in-depth:conformations}{对位交叉构象}，故邻位占三分之二。
\end{solution}

\begin{solution}{exo:b1:stereochemistry-in-depth:6}
碳链竖直，\ce{COOH} 在上，\ce{CH3} 在下；优先次序
$\ce{NH2} > \ce{COOH} > \ce{CH3} > \ce{H}$。将 \ce{NH2} 放在左侧，
H 放在右侧，$\ce{NH2} \to \ce{COOH} \to \ce{CH3}$ 所见为顺时针
；H 水平指向读者，故\omterm{def:b1:stereochemistry-in-depth:isomers}{构型}标记应反转
为 \cip{S}。氨基位于左侧。
\end{solution}

\begin{solution}{exo:b1:stereochemistry-in-depth:7}
\textit{cis}：相邻两碳上一键朝上、一键朝下，于直立、平伏位置中
构成直立–平伏组合，两种\omterm{def:b1:stereochemistry-in-depth:conformations}{椅式构象}中皆然。
\textit{trans}：二平伏或二直立。二平伏\omterm{def:b1:stereochemistry-in-depth:conformations}{椅式构象}中，两个
甲基彼此邻位交叉（一项相互作用）；\textit{cis}
异构体还有一个直立甲基，能量约高 \qty{5.7}{kJ/mol}。
\textit{trans} 异构体较稳定。
\end{solution}

\begin{solution}{exo:b1:stereochemistry-in-depth:8}
$[\alpha] = 13.2/(2.00 \times 0.100) = +66.0^\circ$。若管长 \qty{1.00}{dm}，
则 $+6.60^\circ$。对映体在 \qty{2.00}{dm} 管中为 $-13.2^\circ$。
\end{solution}

\begin{solution}{exo:b1:stereochemistry-in-depth:9}
\qty{25}{\celsius}：$\exp(5700/(8.314 \times 298.15)) = 10.0$，比例
为 $10.0/11.0 = \qty{91}{\percent}$。\qty{-50}{\celsius}：$\exp(5700/(8.314
\times 223.15)) = 21.6$，比例为 \qty{96}{\percent}。
\end{solution}

\begin{solution}{exo:b1:stereochemistry-in-depth:10}
乙烷中六个氢全部等同，各\omterm{def:b1:stereochemistry-in-depth:conformations}{交叉式构象}
皆相同。丁烷中，后端甲基可与前端甲基相对
（\omterm{def:b1:stereochemistry-in-depth:conformations}{对位交叉构象}）或位于其旁（$\pm 60^\circ$，\omterm{def:b1:stereochemistry-in-depth:conformations}{邻位交叉构象}），此时两甲基
互相挤压。布居数：对位为 1，邻位为 $2\exp(-2.8/2.48) = 0.65$，对位约
\qty{61}{\percent}，邻位约 \qty{39}{\percent}。
\end{solution}

\begin{solution}{exo:b1:stereochemistry-in-depth:11}
C2 与 C4 为\omterm{def:b1:stereochemistry-in-depth:stereocentre}{立体中心}。C3 带有 H、OH 及两个
构造相同的两个支链，只有当 C2 与 C4 的
\omterm{def:b1:stereochemistry-in-depth:isomers}{构型}标记相反时两者才有差别，此时 C3 才为\omterm{def:b1:stereochemistry-in-depth:stereocentre}{立体中心}（假不对称碳原子）。异构体：
\cip{2R,4R} 与 \cip{2S,4S}，为一对\omterm{def:b1:stereochemistry-in-depth:enantiomers}{对映异构体}（C3 非\omterm{def:b1:stereochemistry-in-depth:stereocentre}{立体中心}）；
\cip{2R,4S} 因 C3 的两种排列，得两个\omterm{def:b1:stereochemistry-in-depth:enantiomers}{内消旋化合物}。共四种。
\end{solution}

\begin{solution}{exo:b1:stereochemistry-in-depth:12}
$[\alpha]_{\text{样品}} = 2.31/(1.00 \times 0.050) = +46.2^\circ$，且
$2x - 1 = 46.2/66.0 = 0.700$：$x = 0.850$，即 \qty{85}{\percent} 为
$(+)$ 对映体，\qty{15}{\percent} 为 $(-)$ 对映体。
\end{solution}

\begin{solution}{pb:b1:stereochemistry-in-depth:1}
\textbf{1.} 环己烷环：C1 带 OH，C2 带异丙基，C5 带
甲基，C1、C2 与 C5 均为\omterm{def:b1:stereochemistry-in-depth:stereocentre}{立体中心}。
\textbf{2.} $2^3 = 8$，即四对\omterm{def:b1:stereochemistry-in-depth:enantiomers}{对映异构体}。
\textbf{3.} 任一\omterm{def:b1:stereochemistry-in-depth:isomers}{构型}均无对称面：三个取代基
互不相同。
\textbf{4.} \ce{OH} $>$ C2 (C,C,H) $>$ C6 (C,H,H) $>$ H.
\textbf{5.} C1（O,C,H）$>$ 异丙基（C,C,H）$>$ C3（C,H,H）$>$ H。
\textbf{6.} C6 与 C4 均为（C,H,H）；再往下，C6 通向 C1（O,C,H），C4 通向 C3
（C,H,H）：C6 $>$ C4 $>$ \ce{CH3} $>$ H。
\textbf{7.} \cip{1S,2R,5S}.
\textbf{8.} \omterm{def:b1:stereochemistry-in-depth:conformations}{椅式构象}中，相邻碳（C1、C2）上的两个\omterm{def:b1:stereochemistry-in-depth:conformations}{平伏键}
一上一下，为 \textit{trans}；环上碳 1 与 3
（C1 与 C5，经 C6）上两\omterm{def:b1:stereochemistry-in-depth:conformations}{平伏键}指向同一方向：
为 \textit{cis}。此即薄荷醇的\omterm{def:b1:stereochemistry-in-depth:isomers}{构型}。
\textbf{9.} 三者皆变为\omterm{def:b1:stereochemistry-in-depth:conformations}{直立键}。
\textbf{10.} 翻转后，仅甲基一项便约增加
\qty{5.7}{kJ/mol}，较大的异丙基更多；OH 另有较小的
代价，直立 OH 与甲基位于碳 1、3 上，并且在
同侧，二者亦相互挤压。总计远超 \qty{12}{kJ/mol}，
即比值超过 $\exp(12000/2479) \approx 130$ 比 1，薄荷醇
几乎全部处于全平伏的\omterm{def:b1:stereochemistry-in-depth:conformations}{椅式构象}中。
\textbf{11.} 为\omterm{def:b1:stereochemistry-in-depth:enantiomers}{非对映异构体}：三个\omterm{def:b1:stereochemistry-in-depth:stereocentre}{手性}碳原子中仅一个不同。
\textbf{12.} 异丙基与甲基均为\omterm{def:b1:stereochemistry-in-depth:conformations}{平伏键}时，OH 为\omterm{def:b1:stereochemistry-in-depth:conformations}{直立键}。
\textbf{13.} \omterm{def:b1:stereochemistry-in-depth:enantiomers}{非对映异构体}的各项性质均不相同（此处指
直立或平伏 OH 的\omterm{def:b1:intermolecular-forces-solvents:hbond}{氢键}）；\omterm{def:b1:stereochemistry-in-depth:enantiomers}{对映异构体}仅
与\omterm{def:b1:stereochemistry-in-depth:stereocentre}{手性}物质作用时才有差异。
\textbf{14.} $[\alpha] = -2.46/(1.00 \times 0.0500) = -49.2^\circ$，处于
手册所列范围之内。
% ledger: alpha:l-menthol-range
\textbf{15.} $+49.2^\circ$; $0$.
\textbf{16.} $-1.97/0.0500 = -39.4^\circ$.
\textbf{17.} $\alpha = \ell c[\alpha]_{(-)}\bigl(x - (1 - x)\bigr) =
\ell c[\alpha]_{(-)}(2x - 1)$.
\textbf{18.} $x = \frac12(1 + \alpha/\alpha_{\text{参考}})$，在 $\ell$ 与 $c$ 相等时。
\textbf{19.} 是的：其他任何\omterm{def:b1:stereochemistry-in-depth:optical-activity}{旋光性}化合物都会另加一项
（\omterm{prop:b1:stereochemistry-in-depth:biot}{毕奥定律}）。在将旋光度读作组成之前，须先用另一种方法
（色谱法）检验化学纯度。
\textbf{20.} $x = \frac12(1 + 1.97/2.46) = 0.900$.
\textbf{21.} $x = \frac12(1 + 2.40/2.46) = 0.988$.
\textbf{22.} 第二个（\qty{98.8}{\percent}）。
\textbf{23.} $-3.94^\circ$；同样的读数误差此时所占
角度的比例较小。
\textbf{24.} \textbf{$(-)$-薄荷醇占 \qty{90.0}{\percent}}（以及
其对映体占 \qty{10.0}{\percent}）。
\end{solution}
